\section{Scalar finite-part boundary}
\label{app:scalar-boundaries}

This appendix records the only scalar-boundary fact needed by the present
finite-group theorem chain.  The scalar finite part
\(\mathcal M(A)=\gamma+\FP(A)\) is the trivial-channel normalisation in
Main Theorems~II--III.  It depends on the base matrix \(A\), while the
non-trivial class data depend on the fixed-label twisted determinants of
\((A,\tau,G)\).  Thus scalar finite parts and the canonical cyclic-lift
records of \Cref{app:cyclic-lift-comparison} are retained here only as
normalisation and comparison tools.

Questions about holomorphic variation of \(\FP(A)\), effective observation of
truncations, arithmetic refinements of root-of-unity records, and
growing-family estimates are not used as hidden inputs for Main
Theorems~I--III.

\begin{lemma}[Scalar records enter only through the trivial channel]
  \label{lem:scalar-boundary-trivial-channel-only}
  Fix a primitive matrix \(A\), a finite group \(G\), and a non-empty family
  \(\mathscr T\subset\mathcal T_{\mathrm{gap}}(A,G)\) of edge cocycles for
  which the hypotheses of Main Theorems~II--III hold.  Let \(\mathcal R(A)\)
  be any scalar fixed-matrix record depending only on \(A\), for instance
  \(\FP(A)\), \(\mathcal M(A)\), the reduced determinant \(C(A)\), or a
  canonical cyclic-lift record named in \Cref{app:cyclic-lift-comparison}.
  Then \(\mathcal R(A)\) is constant as \(\tau\in\mathscr T\) varies and can
  determine the normalised class-Mertens profile of Main Theorem~II on
  \(\mathscr T\) only up to the trivial projection
  \(\mathcal M(A)\mathbf1_G\).  It determines the non-trivial class profile
  only in the degenerate case where all non-trivial Perron-point channels
  \(\Pi_\varrho^\tau(\lambda^{-1})\), \(\varrho\ne\mathbf1\), are already
  independent of \(\tau\) on \(\mathscr T\).
\end{lemma}

\begin{proof}
  The fixed-matrix record \(\mathcal R(A)\) does not involve the cocycle
  \(\tau\), so it is constant on \(\mathscr T\).  By
  \Cref{cor:fixed-matrix-scalar-interface-only}, the constant projection of
  the normalised class-Mertens profile is exactly
  \(\mathcal M(A)\mathbf1_G\), and the complementary zero-mean projection is
  the class function with Peter--Weyl coefficients
  \(\Pi_\varrho^\tau(\lambda^{-1})\) for
  \(\varrho\ne\mathbf1\).  A scalar record depending only on \(A\) can
  recover that complementary projection throughout \(\mathscr T\) only if the
  complementary projection itself is constant on \(\mathscr T\), equivalently
  only if all displayed non-trivial coefficients are independent of \(\tau\).
\end{proof}

Consequently the finite-group information in Main Theorems~I--III comes from
the fixed-label determinant family and the Adams--M\"obius inversion in the
main text, not from scalar finite-matrix records alone.
